\section{MOTIVATION}
\label{sec:motivation}

``What kind and how much information is needed to achieve a desired type of control?'' \cite{kalman_1959}. This fundamental question, posed by Rudolf E. Kalman over six decades ago, remains surprisingly current. While Kalman was addressing the mathematical foundations of control systems, establishing that a system is observable only if its internal states can be inferred from its external outputs, the software industry today faces an analogous dilemma in a vastly more complex environment: the cloud-native ecosystem.

In the era of distributed systems, the ``plant'' — to use control theory terminology — has evolved from a linear, stationary entity into dynamic, ephemeral, and distributed clusters orchestrated by Kubernetes. Just as Kalman established the duality principle, stating that controllability and observability are two sides of the same coin, modern software engineering is discovering that it is impossible to effectively ``control'' (orchestrate, scale, and heal) complex microservices without deep ``observability''.

However, implementing this concept in Kubernetes environments is far from trivial. The transition from monolithic architectures to microservices has exploded the cardinality of data and the complexity of interactions. The answer to ``How to implement observability?'' is often as ambiguous as the architectural decisions explored in previous studies: ``It depends.''

To bridge the conceptual gap between classical control theory and modern cloud-native ecosystems, Kalman's principles must be adapted into an \textit{architectural metaphor} rather than a strict mathematical proof. In the context of this research, the Kubernetes cluster acts as the ``plant'' (the dynamic system being controlled). The ``internal states'' of this system encompass the functional health of microservices, memory allocation, network latency between nodes, and queue saturation. The ``external outputs'' are the telemetry signals (metrics, logs, and traces) emitted by the workloads. Therefore, achieving observability in Kubernetes does not strictly mean calculating algebraic state matrices; rather, it implies designing a telemetry architecture that structurally reduces ambiguity in state identification. By semantically correlating these outputs, the architecture empowers the Control Plane—or a human Site Reliability Engineer (SRE)—to apply corrective feedback with deterministic precision.

The Cloud Native Computing Foundation (CNCF) landscape offers a plethora of open-source tools—Prometheus, OpenTelemetry, Jaeger or Grafana Tempo, Fluentd, among others. Yet, the sheer volume of choices creates a paradox of choice. Organizations struggle not with a lack of data, but with a lack of cohesive strategy. They often assemble fragmented pipelines that generate noise rather than insight, failing to achieve the actionable visibility required to regulate the system.

Current literature and industry practices highlight several friction points in this domain:

\begin{itemize}
    \item \textbf{Tool Sprawl and Fragmentation}: The challenge of integrating disparate open-source tools into a unified pane of glass.
    \item \textbf{Correlation Context}: The difficulty in linking metrics, logs, and traces to understand causality in distributed failures.
    \item \textbf{Performance Overhead}: The balance between granular data collection and the resource cost on the cluster.
    \item \textbf{Cognitive Load}: The human operator's ability to interpret vast amounts of telemetry data to make architectural decisions.
\end{itemize}

Much like the need for documented Architectural Decisions (ADs) to prevent knowledge decay in software design, there is a pressing need for a structured approach to observability infrastructure. It is not enough to simply deploy disconnected agents. Instead, observability must be engineered as a closed feedback loop: \textit{telemetry} (signal generation and collection) drives \textit{diagnosis} (causal correlation), which in turn informs \textit{orchestration} (actionable remediation). By framing observability through this control-theory metaphor, Site Reliability Engineering practices shift from reactive troubleshooting to proactive, data-driven system regulation.

Therefore, the motivation for this thesis stems from the gap between the theoretical necessity of observability for system control and the practical lack of standardized, open-source reference architectures for Kubernetes. By revisiting these foundational concepts where observability is a prerequisite for optimal regulation, this work seeks to move beyond ad-hoc implementations. The goal is to propose a reference architecture that serves as a blueprint for building robust, observable cloud-native systems, enabling practitioners to transform raw telemetry into actionable control signals.